\documentclass[conference]{IEEEtran}
\usepackage[utf8]{inputenc}
\usepackage[T1]{fontenc}
\usepackage{cite}
\usepackage{amsmath,amssymb}
\usepackage{graphicx}
\usepackage{booktabs}
\usepackage{url}

\title{Generate--Verify--Repair for Intent‑Driven Network Changes: Controlled Evaluation with a Deterministic Verifier}

\author{
\IEEEauthorblockN{Autonomous AI Researcher}
\IEEEauthorblockA{CNSM 2026 Agentic AI Researcher Track}
}

\begin{document}

\maketitle

\begin{abstract}
We evaluate whether a generate--verify--repair (GVR) pipeline---single‑shot LLM generation followed by a deterministic network verifier and up to one automated repair---reduces accepted unsafe network changes versus LLM‑only generation on a standardized synthetic NetOps intent workload. In a preregistered paired design (study-hosted-netops-gvr-v1-2026-08-29) we ran N=500 distinct synthetic intents (shared single initial candidate per intent) with generation by gpt-5-mini and adjudication by a deterministic verifier. Results: guarded (GVR) accepted 500/500 final changes; baseline (LLM-only) accepted 496/500 (paired difference = 0.008). The preregistered primary exact conditional McNemar two‑sided test on the discordant pairs (n10 = 4, n01 = 0) yields p = 0.125 (computed as p = 2 * PBinom(X >= 4; n = 4, p = 0.5)). An exploratory paired bootstrap percentile 95\% CI (10,000 resamples, seed = 1729) = [0.002, 0.016] (bootstrap labeled exploratory per preregistration and interpreted cautiously given only four discordants). Repairs were invoked for 4 initial candidates; total model calls used = 504 (500 shared\_initial + 4 repair calls). The preregistered templated‑automation comparator was not executed. Because the preregistered decision rule is the exact conditional McNemar (p = 0.125), we do not claim confirmatory statistical significance. We reconcile the conditional McNemar and the exploratory bootstrap CI, document verifier scope and known blind spots, provide execution accounting and representative validator traces, and give explicit artifact pointers and hashes for reproducibility.
\end{abstract}

\section{Introduction}
Operator intent\ensuremath{\rightarrow}device configuration translation can reduce manual effort, but errors in generated CLI sequences or transient unsafe states can cause outages. Modern large language models (LLMs) show promise for translating human intent to device configuration, yet hallucination and factual errors remain central reliability concerns for operational NetOps. A practical mitigation architecture is generate--verify--repair (GVR): produce a candidate configuration with an LLM, deterministically verify it against encoded workflow and policy rules, and, if verification fails, run a bounded automated repair attempt. This paper reports a preregistered, artifact‑grounded paired experiment that asks: Does a GVR pipeline (LLM \ensuremath{\rightarrow} deterministic verifier \ensuremath{\rightarrow} up to one automated repair) produce fewer accepted unsafe changes than LLM‑only generation on a standardized synthetic NetOps intent workload? We contribute: (1) a preregistered paired trial (N=500 paired intents) executed end‑to‑end with archived artifacts; (2) an artifact‑grounded description of generation, verifier semantics, and the bounded repair loop; (3) reconciled confirmatory inference (two‑sided conditional exact McNemar) and exploratory paired bootstrap CI descriptions with an explicit small‑discordant‑count caveat; and (4) execution accounting, representative validator traces, and a discussion of verifier blind spots and practical external‑validity limits.

\section{Related Work}
Three converging literatures motivate and contextualize this study. First, recent surveys and system papers document active interest in applying LLMs to NetOps tasks such as intent translation, configuration generation, AIOps, and multi‑agent orchestration. These works show practical system designs and report deployment experiences or prototypes while consistently highlighting reliability risks from hallucination and factual errors in generated artifacts [doi:10.1145/3718958.3750537; https://openalex.org/W4415071524; doi:10.1002/nem.2313]. They motivate architectures that pair generative models with domain tooling rather than relying on unconstrained generation alone.

Second, methodological literature on tool‑augmented LLMs and generate‑validate‑repair pipelines argues that external deterministic or domain tools (verifiers, simulators, parsers) materially change the failure‑mode profile of LLM outputs and therefore should be central to evaluation design. Several studies propose or evaluate RAG/tool‑augmented workflows and stress the need for standardized evaluation practices that explicitly measure verifier coverage, oracle mismatch, and tool‑call cost tradeoffs [https://openalex.org/W4403443637; doi:10.2139/ssrn.6647800]. This body of work clarifies two practical points relevant to our design: (1) verification tools provide high‑precision adjudication for well‑specified properties (syntactic correctness, required sequencing, availability constraints) but rarely capture all production hazards (vendor idiosyncrasies, timing/convergence effects); and (2) the operational cost of invoking heavyweight verifiers or iterative repair loops must be measured alongside safety benefits.

Third, the verification and network‑analysis literature supplies deterministic or formal checkers capable of validating many configuration properties at scale. Recent verifier and distributed‑verification systems adapt Batfish‑style analyses and localized sub‑specification approaches to large inventories, enabling automated checks of reachability, required sequences, and group/availability constraints [https://openalex.org/W4413757123; doi:10.1145/3718958.3750516]. These verifiers form the practical oracle in generate--verify workflows but, as prior work emphasizes, their utility depends on the modeled properties and the abstraction fidelity used to simulate device behavior. Our study builds on these foundations by pairing a compact deterministic verifier (explicitly encoding required‑sequence, group/availability, paired‑field, and final‑state checks) with a tightly bounded repair policy and measuring end‑to‑end outcomes using preregistered paired inference.

Gaps and relation to the present experiment. While prior systems integrate verification into LLM‑driven NetOps frameworks, few published controlled, preregistered experiments compare end‑to‑end GVR pipelines against LLM‑only or templated automation baselines at the paired‑intent scale with full artifact archival. The methodological calls for standardized evaluation (tool budgets, verifier coverage reporting, and paired designs) directly informed our preregistration. Our contribution is therefore empirical and procedural: a preregistered, paired measurement using a deterministic verifier as the adjudicating oracle and an explicitly bounded single‑repair policy, with complete archival of generation outputs, verifier traces, and analysis artifacts to permit independent reconstruction and targeted follow‑up studies that address verifier fidelity and production realism.

\section{Methodology}
Preregistration and confirmatory plan: The study was preregistered as study-hosted-netops-gvr-v1-2026-08-29. The primary estimand is the paired success‑rate difference (GVR - baseline) across N = 500 distinct intents stratified evenly across 10 synthetic topology profiles (50 intents per profile). Sampling seeds, the deterministic task generator, and the analysis executor were frozen before execution (preregistration SHA‑256 = aa8982a2\allowbreak{}7c73e515\allowbreak{}21ac023a\allowbreak{}78adcfa3\allowbreak{}58ef3978\allowbreak{}c0f9bc05\allowbreak{}213801fc\allowbreak{}bd2f1d0a). The preregistered confirmatory test is the two‑sided conditional exact McNemar on discordant pairs; an exploratory paired bootstrap percentile 95\% CI (10,000 resamples; bootstrap\_seed = 1729) was pre‑specified as a descriptive interval.

Generation, orchestration, and empirical call accounting: For each intent a single shared\_initial\_generation candidate was produced by the declared model gpt-5-mini (declared\_model\_version recorded in execution/model\_configuration.json; execution/model\_configuration.json SHA‑256 = a40e96e2\allowbreak{}12ab87da\allowbreak{}251ac0d6\allowbreak{}dbb75bc5\allowbreak{}43afa134\allowbreak{}38b5a6b9\allowbreak{}ebd07359\allowbreak{}7ffdd820). The GVR controller validated the shared candidate with the deterministic verifier deterministic\_netops\_validator\_v5. If verification failed, the controller invoked at most one automated repair attempt (maximum\_repair\_calls\_per\_task = 1). Provider‑call accounting recorded hosted\_model\_call\_event\_count = 504 (stage\_counts: shared\_initial\_generation = 500; repair = 4) and model\_calls\_used = 504 in the execution manifest; these counts are part of the archived provider audit artifacts.

Deterministic verifier semantics, emitted fields, and artifact pointer: The primary verifier implements a deterministic validation pipeline on a normalized CLI sequence comprising the following stages: (1) syntactic parsing and canonical normalization of CLI lines; (2) workflow and required‑sequence checks (patterns such as must\_be\_down\_for\_fields, all\_down\_before\_any\_field\_change); (3) availability/group constraints (exactly\_active, minimum\_active\_in\_groups); (4) dependency and paired‑field enforcement (paired MTU changes, equal\_final\_fields); and (5) final‑state comparison with the task expected\_final\_state. The verifier emits, per episode, a human‑readable validator\_trace including normalized\_configuration, observed\_touched\_field\_count, parsed\_command\_count, required\_command\_count, violation list and codes, validator\_id/version, and a boolean valid flag. The human‑readable verifier codebook (violation‑code \ensuremath{\rightarrow} short description) is enumerated in the archived execution manifest; the execution manifest file (execution/execution\_manifest.json) in the evidence bundle lists the exact verifier codebook artifact path and its artifact hash for direct retrieval. Representative verifier\_trace JSON objects are archived under execution/scoring/ and include the emitted fields described above.

Repair policy and repair‑prompt semantics: If the verifier fails a shared\_initial candidate, the GVR controller issues at most one repair prompt to the same generation model using a constrained repair template that includes: (a) the original intent abstraction; (b) the verifier violation codes and short descriptions; and (c) a request for a corrected canonical CLI sequence. The repair template is intentionally minimal to limit model‑call overhead and to isolate the marginal effect of a single repair attempt. A repair candidate is accepted only if the verifier returns valid=true for the subsequent candidate. Repair invocations and their acceptance are logged per episode in execution/scoring/ and in the provider traces; in this run 4 repair calls were invoked and all resulted in verifier passes.

Statistical procedures, exact McNemar implementation, and bootstrap specification: The preregistered primary decision rule is the two‑sided conditional exact McNemar computed on discordant pairs only. Let n10 be the count of pairs where guarded succeeds and baseline fails, and n01 the reverse; let n = n10 + n01 and k = n10. The exact two‑sided p‑value used in the preregistration is computed as p = 2 * PBinom(X >= k; n, 0.5) with clipping at 1.0 when required. This implementation conditions on the observed margin n and tests whether the probability of observing at least k successes under X \textasciitilde{} Binomial(n, 0.5) is small. The preregistered exploratory interval procedure is a paired bootstrap percentile 95\% CI: resample the N paired outcomes with replacement 10,000 times (bootstrap\_resamples = 10,000; bootstrap\_seed = 1729), compute the paired success‑rate difference per resample, and take the 2.5th and 97.5th percentiles. The bootstrap was explicitly predeclared exploratory in the analysis plan. Because the conditional McNemar conditions on the observed discordant count n and evaluates the exact binomial tail, whereas the bootstrap resamples pairs without fixing margins, the two procedures can diverge when margins are nearly fixed and discordant counts are small; this divergence and its implications are addressed in the Discussion.

Reproducibility and logged artifacts referenced from methodology: The execution manifest and associated artifact hash manifest contain all paths and artifact hashes referenced above (representative entries include execution/raw\_results.jsonl SHA‑256 = fe9b93ef\allowbreak{}5cc5a237\allowbreak{}a3361f18\allowbreak{}f45fabf1\allowbreak{}4cd23a1c\allowbreak{}679556f6\allowbreak{}5477f9f6\allowbreak{}81915d02; execution/task\_manifest.jsonl SHA‑256 = 5a822e78\allowbreak{}7302a0b3\allowbreak{}bb03a86a\allowbreak{}81b20564\allowbreak{}a44e2636\allowbreak{}731f853f\allowbreak{}9d45ac12\allowbreak{}e4876cc8; analysis/paired\_contingency\_table.csv SHA‑256 = 8fb135ce\allowbreak{}c541cbe0\allowbreak{}b14f1612\allowbreak{}5db33154\allowbreak{}a41460c4\allowbreak{}9958427f\allowbreak{}e8e7f3e1\allowbreak{}f2abfece). The execution manifest (execution/execution\_manifest.json) enumerates the verifier codebook artifact path and its artifact hash for direct retrieval from the archived bundle. Where specific numeric analysis parameters were pre‑specified, they are reported here: bootstrap\_seed = 1729 and bootstrap\_resamples = 10,000. The preregistration SHA and model\_configuration.json SHA are recorded above and in the archived manifest.

\section{Results}
Execution and primary counts: The run completed completed\_episode\_count = 1000 (500 intents \ensuremath{\times} 2 conditions) with model\_calls\_used = 504. analysis/condition\_summary.csv (archived hash: f3e19742\allowbreak{}5cc03dec\allowbreak{}41cda77f\allowbreak{}078f6d0b\allowbreak{}079b06bb\allowbreak{}de7ac736\allowbreak{}de5fe988\allowbreak{}b027ac58) reports baseline successes 496/500 (0.992) and guarded successes 500/500 (1.000). The paired contingency table (analysis/paired\_contingency\_table.csv; archived hash: 8fb135ce\allowbreak{}c541cbe0\allowbreak{}b14f1612\allowbreak{}5db33154\allowbreak{}a41460c4\allowbreak{}9958427f\allowbreak{}e8e7f3e1\allowbreak{}f2abfece) is: n11 = 496, n10 = 4, n01 = 0, n00 = 0. Repairs were invoked for four initial candidates; all four repair attempts resulted in subsequent verifier passes and correspond to the four discordant pairs (n10 = 4). Mean model calls per accepted change \ensuremath{\approx} 504/500 = 1.008.

Confirmatory exact McNemar calculation: Observed discordants n = n10 + n01 = 4 with k = n10 = 4. The preregistered conditional McNemar two‑sided p‑value is p = 2 * PBinom(X >= 4; n = 4, p = 0.5). For n = 4, PBinom(X >= 4; n = 4, p = 0.5) = 1/16, so p = 2 * (1/16) = 0.125. Because this exact conditional test was the preregistered decision rule, we do not claim confirmatory statistical significance for the primary estimand.

Exploratory paired bootstrap interval: The preregistered exploratory paired bootstrap percentile 95\% CI (10,000 resamples; seed = 1729; bootstrap artifacts archived as analysis/paired\_difference\_figure.svg and analysis/analysis\_log.jsonl; analysis\_log.jsonl SHA‑256 = 261ec80e\allowbreak{}a9343fcc\allowbreak{}bcd0a622\allowbreak{}94f53827\allowbreak{}73770fd5\allowbreak{}9d387dc9\allowbreak{}d749de50\allowbreak{}ca909b3e) for the paired difference is [0.002, 0.016]. The bootstrap was pre‑specified as exploratory; as discussed below, the bootstrap can produce a CI that excludes zero even when the conditional McNemar is non‑significant when margins are nearly fixed and discordant counts are small.

Representative artifact diagnostics: Representative scoring JSON objects (execution/scoring/task-00000*-*.json) show verifier\_trace entries with validation\_before, validation\_after, normalized\_configuration, parsed\_command\_count, required\_command\_count, violation\_count, and validator\_version = 5.0. Example episodes covering medium and hard difficulty patterns (shutdown\_configure\_restore, make\_before\_break\_uplink\_switchover, paired\_link\_atomic\_mtu\_change) are archived and referenced in the representative\_tasks list of the evidence bundle. The four repaired episodes show verifier‑detectable sequence or availability violations that a single constrained repair prompt corrected; repair prompts included the verifier violation codes and the intent abstraction per the preregistered policy.

Provider audit and execution manifest: Deterministic\_reconciliation.json records provider\_call\_audit: hosted\_model\_call\_event\_count = 504, completed\_hosted\_model\_call\_event\_count = 504, failed\_hosted\_model\_call\_event\_count = 0, cache\_hit\_count = 0, cache\_miss\_count = 504, and stage\_counts: shared\_initial\_generation = 500 and repair = 4. Representative raw artifacts and hashes cited here include execution/raw\_results.jsonl (SHA‑256 = fe9b93ef\allowbreak{}5cc5a237\allowbreak{}a3361f18\allowbreak{}f45fabf1\allowbreak{}4cd23a1c\allowbreak{}679556f6\allowbreak{}5477f9f6\allowbreak{}81915d02) and execution/task\_manifest.jsonl (SHA‑256 = 5a822e78\allowbreak{}7302a0b3\allowbreak{}bb03a86a\allowbreak{}81b20564\allowbreak{}a44e2636\allowbreak{}731f853f\allowbreak{}9d45ac12\allowbreak{}e4876cc8). The paired contingency CSV hash is analysis/paired\_contingency\_table.csv (SHA‑256 = 8fb135ce\allowbreak{}c541cbe0\allowbreak{}b14f1612\allowbreak{}5db33154\allowbreak{}a41460c4\allowbreak{}9958427f\allowbreak{}e8e7f3e1\allowbreak{}f2abfece). The preregistration artifact and master prompt are archived (preregistration SHA‑256 = aa8982a2\allowbreak{}7c73e515\allowbreak{}21ac023a\allowbreak{}78adcfa3\allowbreak{}58ef3978\allowbreak{}c0f9bc05\allowbreak{}213801fc\allowbreak{}bd2f1d0a; initial master prompt path provenance/master\_prompt.txt SHA‑256 = 1872df1e\allowbreak{}1805d2d9\allowbreak{}6940456c\allowbreak{}a016bd66\allowbreak{}5d1d5196\allowbreak{}add77f5a\allowbreak{}cdf1582b\allowbreak{}b39b15ba). All artifacts cited are retrievable from the archived execution manifest in the evidence bundle.

\section{Discussion}
Interpretation and reconciliation of inferential outcomes: The GVR pipeline prevented four verifier‑detectable unsafe initial candidates that baseline would have accepted. Under the preregistered conditional exact McNemar (p = 0.125) this difference is not confirmatory. The exploratory bootstrap CI excludes zero but should be interpreted cautiously: the conditional McNemar conditions on observed margins and evaluates the binomial tail of discordant outcomes; with n = 4 and k = 4 this yields p = 0.125. The bootstrap resamples pairs without fixing margins and can produce a nonzero lower CI bound when most pairs are concordant and a small number are discordant. Because the bootstrap was predeclared exploratory and because small discordant counts reduce frequentist coverage guarantees for percentile intervals, we report both inferential perspectives and follow the preregistered decision rule for the primary claim.

Verifier coverage and concrete blind spots: The deterministic verifier enforces canonical normalization, required sequencing, availability/group constraints, dependency and paired‑field rules, and exact final‑state equality to the task expected\_final\_state. It does not model vendor‑specific CLI semantics, timing/convergence effects in control‑plane protocols, vendor idiosyncrasies, hardware resource constraints, or out‑of‑band side effects. Passing this verifier is necessary for experiment safety but not sufficient for all production hazards; expanding verifier fidelity (vendor parsers, control‑plane simulators, timing models) is required before operational automation. The human‑readable verifier codebook (violation‑code \ensuremath{\rightarrow} description) is enumerated in the execution manifest (see execution/execution\_manifest.json in the evidence bundle) and the manifest lists the exact codebook artifact path and its SHA‑256 for direct retrieval.

Operational implications and cost tradeoffs: A compact deterministic verifier plus a bounded repair loop intercepted verifier‑detectable hazards at low LLM overhead in this synthetic workload. Observed repair rate was 4/500 (0.8\%), implying modest additional model calls under the bounded policy. Production adoption requires (a) measuring verifier runtime and scalability on real inventories, (b) extending verifier fidelity to vendor/timing effects, and (c) evaluating adversarial inputs where repair frequency and cost may increase. The provider audit (hosted\_model\_call\_event\_count = 504) quantifies the empirical model‑call overhead for this run.

Threats to validity and generality: Internal validity is strong because of preregistration, fixed seeds, and complete pairs (complete\_pair\_count = 500; execution\_manifest failed\_episode\_count = 0). Construct validity is limited by verifier abstractions and synthetic device templates (synthetic\_topologies\_v1 and synthetic\_device\_configs\_v1). External validity is constrained by single‑model scope (gpt-5-mini) and synthetic topologies; multi‑model replication and real vendor CLI tests are required to generalize. Conclusion validity was affected by realized margins: the preregistered power plan targeted an absolute paired increase \ensuremath{\approx} 0.08; the observed near‑ceiling baseline (0.992) produced only four discordant pairs and reduced realized power relative to the preregistered assumption. These limitations are documented in the archived manifest and preregistration.

\section{Limitations}
\begin{itemize}
\item Synthetic tasks and topology profiles were used (synthetic\_topologies\_v1 and synthetic\_device\_configs\_v1); external validity to production hardware, vendor‑specific CLIs, live telemetry, and control‑plane timing effects is untested.
\item Only a single generation model (gpt-5-mini) was evaluated; multi‑model generality and replication remain to be established.
\item Safety in this experiment is defined relative to deterministic\_netops\_validator\_v5's encoded rule set (syntactic parsing/normalization, required‑sequence/workflow‑policy checks, protected‑interface rules, group and availability constraints, dependency‑rule enforcement, and final‑state comparison to the task expected\_final\_state). Passing this verifier is necessary for the experiment's outcome but not sufficient for all real‑world safety properties.
\item The preregistered power plan targeted an absolute paired increase \ensuremath{\approx} 0.08 (8 percentage points). The observed baseline success rate (496/500 = 0.992) produced only four discordant pairs, substantially reducing realized power relative to the preregistered assumption and limiting confirmatory sensitivity.
\item The preregistered templated‑automation comparator (secondary confirmatory arm) was not executed; therefore the preregistered multiplicity plan (two confirmatory tests with Bonferroni correction) could not be fully applied and the familywise error interpretation is partial.
\end{itemize}

\section{Conclusion}
In a preregistered paired trial (N = 500) using gpt-5-mini and a deterministic verifier, a generate--verify--repair pipeline repaired four verifier‑detectable unsafe initial candidates at modest model‑call cost (4 repairs; model\_calls\_used = 504). The paired difference = 0.008 yields an exploratory bootstrap 95\% CI [0.002, 0.016], while the preregistered exact conditional McNemar test (n10 = 4, n01 = 0) yields p = 0.125; under the preregistered decision rule we do not claim confirmatory significance. Deterministic verification plus a bounded repair loop can remove verifier‑detectable unsafe outputs with low overhead in synthetic settings; broader operational claims require expanded verifier fidelity, adversarial NetOps evaluation, multi‑model replication, and execution of the preregistered templated comparator (which was not executed in this run). All primary execution and analysis artifacts cited here are archived in the evidence bundle for independent verification; see the archived execution manifest for full artifact paths and hashes.

\section*{Disclosure Statement}
This study was executed autonomously after an explicit autonomy lock. Generation model used in the archived run: gpt-5-mini (declared\_model\_version recorded in execution/model\_configuration.json; execution/model\_configuration.json SHA‑256 = a40e96e2\allowbreak{}12ab87da\allowbreak{}251ac0d6\allowbreak{}dbb75bc5\allowbreak{}43afa134\allowbreak{}38b5a6b9\allowbreak{}ebd07359\allowbreak{}7ffdd820). Orchestration adapter family: hosted\_netops\_gvr\_v1 using the hosted provider recorded as 'openai\_responses'. Key automated components recorded in the run: deterministic\_netops\_task\_generator\_v5 (task synthesis), deterministic\_netops\_validator\_v5 (primary deterministic verifier/scorer), and the GVR controller implementing shared\_initial\_generation plus an automated single‑repair step (maximum\_repair\_calls\_per\_task = 1). Preregistration: study-hosted-netops-gvr-v1-2026-08-29 (preregistration SHA‑256 = aa8982a2\allowbreak{}7c73e515\allowbreak{}21ac023a\allowbreak{}78adcfa3\allowbreak{}58ef3978\allowbreak{}c0f9bc05\allowbreak{}213801fc\allowbreak{}bd2f1d0a). The immutable initial master prompt is archived at provenance/master\_prompt.txt (SHA‑256 = 1872df1e\allowbreak{}1805d2d9\allowbreak{}6940456c\allowbreak{}a016bd66\allowbreak{}5d1d5196\allowbreak{}add77f5a\allowbreak{}cdf1582b\allowbreak{}b39b15ba). Primary high‑level execution quantities reported here (N = 500 complete pairs; model\_calls\_used = 504; repair calls = 4; paired counts n11 = 496, n10 = 4, n01 = 0, n00 = 0) are derived from archived execution and analysis artifacts. Representative artifact hashes included in the evidence bundle: execution/raw\_results.jsonl SHA‑256 = fe9b93ef\allowbreak{}5cc5a237\allowbreak{}a3361f18\allowbreak{}f45fabf1\allowbreak{}4cd23a1c\allowbreak{}679556f6\allowbreak{}5477f9f6\allowbreak{}81915d02; execution/task\_manifest.jsonl SHA‑256 = 5a822e78\allowbreak{}7302a0b3\allowbreak{}bb03a86a\allowbreak{}81b20564\allowbreak{}a44e2636\allowbreak{}731f853f\allowbreak{}9d45ac12\allowbreak{}e4876cc8; analysis/paired\_contingency\_table.csv SHA‑256 = 8fb135ce\allowbreak{}c541cbe0\allowbreak{}b14f1612\allowbreak{}5db33154\allowbreak{}a41460c4\allowbreak{}9958427f\allowbreak{}e8e7f3e1\allowbreak{}f2abfece; analysis/analysis\_log.jsonl SHA‑256 = 261ec80e\allowbreak{}a9343fcc\allowbreak{}bcd0a622\allowbreak{}94f53827\allowbreak{}73770fd5\allowbreak{}9d387dc9\allowbreak{}d749de50\allowbreak{}ca909b3e. The human‑readable verifier codebook (violation‑code \ensuremath{\rightarrow} description) and the full list of artifact paths and hashes are enumerated in the archived execution manifest (execution/execution\_manifest.json); the execution manifest lists the exact verifier codebook artifact path and its SHA‑256 for direct retrieval. The design, preregistration, code generation, execution, analysis, manuscript drafting, autonomous peer review, and revision were performed autonomously after the autonomy lock; no human manually edited or rewrote manuscript technical content after that lock. The preregistered templated‑automation comparator was not executed in this archived run; that unexecuted arm means the originally preregistered two‑test Bonferroni familywise plan could not be fully applied.

\begin{thebibliography}{99}
\bibitem{ref1} doi:10.1145/3718958.3750537
\bibitem{ref2} Sebastian Simon, Alina Mailach, Johannes Dorn, Norbert Siegmund, "A Methodology for Evaluating RAG Systems: A Case Study On Configuration Dependency Validation", 2024, 10.48550/arxiv.2410.08801
\bibitem{ref3} Nguyen Tu, Sukhyun Nam, James Won‐Ki Hong, "Intent‐Based Network Configuration Using Large Language Models", 2024, 10.1002/nem.2313
\bibitem{ref4} Lingzhe Zhang, Tong Jia, Mengxi Jia, Yifan Wu, Aiwei Liu, Yong Yang, Zhonghai Wu, Xuming Hu, Philip S. Yu, Ying Li, "A Survey of AIOps in the Era of Large Language Models", 2025, 10.1145/3746635
\bibitem{ref5} Zhaodong Wang, Samuel Lin, Guanqing Yan, Soudeh Ghorbani, Minlan Yu, Jiawei Zhou, Nathan Hu, Lopa Baruah, Sam Peters, Srikanth Kamath, Jian Yang, Ying Zhang, "Intent-Driven Network Management with Multi-Agent LLMs: The Confucius Framework", 2025, 10.1145/3718958.3750537
\end{thebibliography}

\end{document}
